\documentclass[10pt,a4paper]{amsart}
\usepackage[margin=19mm]{geometry}
\usepackage{amsmath,amssymb,mathtools}
\usepackage[scr=rsfs,cal=euler]{mathalfa}
\usepackage{xfrac}
\usepackage[unicode,hidelinks,bookmarks=false]{hyperref}
\newcommand{\ie}{i.e.\ }
\newcommand{\cf}{cf.\ }
\newcommand{\resp}{respectively\ }
\newcommand{\Subsection}{\subsection}
\providecommand{\nobreakdash}{\nobreak\hbox{-}}
\setlength{\emergencystretch}{2em}
\multlinegap0pt
\usepackage{amssymb}
\usepackage{xcolor}
\usepackage{enumerate}
\usepackage{bm}
\usepackage{dsfont}
\usepackage{mathtools}
\newtheorem{lemma}{Lemma}
\newtheorem{theorem}{Theorem}
\newcommand{\N}{\mathbb{N}}
\newcommand{\Sone}[1]{\Stsigma/\Stsigma (t^{\de}-#1)} %{\mathbf{S}_{t^m-#1}}
\newcommand{\Stsigma}{S[t;\sigma]}
\newcommand{\de}{\mathbf{n}}  %temporary symbol for the degree of $f$
\newcommand{\ring}{S}
\newcommand{\scz}{\mathsf{z}}  %in the first proposition
\title{When isometry and equivalence for skew constacyclic codes coincide}
\author{Monica Nevins, Susanne Pumpluen}
\date{Source: arXiv:2508.06695v3 (CC BY 4.0)}
\begin{document}
\maketitle

\noindent\emph{Source and licence.} When isometry and equivalence for skew constacyclic codes coincide. Version arXiv:2508.06695v3 (CC BY 4.0), distributed by arXiv under the Creative Commons Attribution 4.0 International licence, http://creativecommons.org/licenses/by/4.0/, as recorded in the arXiv licence metadata for this exact version and shown on the abstract page https://arxiv.org/abs/2508.06695v3. Full licence notice: \texttt{licenses/arxiv-2508.06695-CC-BY-4.0.txt}.\par\smallskip

\noindent\textit{Editorial scope.} Counted unit: the article's theorem T:powerassociative (source0.tex lines 535--551). The article's complete original proof of it is delivered with it (source0.tex lines 554--658, one \texttt{proof} environment of 105 lines). No mathematics is written, repaired or re-derived: the statement and the proof are the article's own lines, in the article's own notation, and no retained line is substituted or re-flowed.  Retained uncounted as the source-local context the proof needs: lines 481--497.  Nothing was omitted from any retained span, so the package carries no omission record.  No retained line contains a citation, so the wrapper carries no bibliography and no bibliography entry.  No retained line contains a figure environment, an image-inclusion command or drawing code, so no image is delivered and the package declares no asset dependency.  Editorial formatting only, in addition: an AMS \texttt{amsart} preamble that restates the article's own declarations and macros used by the retained lines, namely the environments \texttt{lemma}, \texttt{theorem} on separate counters, together with the standard packages those lines need.  The retained environments print in this document as Lemma 1, Theorem 1 in order of appearance, whereas the article prints the same items with the numbering of its own full document.  The complete native source of the article as distributed in the arXiv e-print for this exact version is bundled unchanged as \texttt{sources/182958/source0.tex}.
\begin{lemma}\label{L:keystep}
    Suppose $\de,k\in \mathbb{N}$ and $\alpha,b\in \ring$ are not zero divisors.  Suppose
\begin{equation}\label{E:conditiononassociativity}
    \alpha b = \sigma^{\de}(\alpha) \sigma^k(b).
    \end{equation}
    Then for all $s,\ell\in \mathbb{N}$ we have
    \begin{equation}\label{E:keyequality}
\sigma^{[sk]_{\de}}
(N_{\ell}^{\sigma^k}(\alpha))
\prod_{j=1}^z\sigma^{sk-j{\de}}(b) =
\sigma^{sk}
(N_{\ell}^{\sigma^k}(\alpha))
\prod_{j=1}^z\sigma^{sk-j{\de}}
(\sigma^{\ell k}(b)),
    \end{equation}
    where $z= \lfloor \frac{sk}{{\de}}\rfloor$ and $[sk]_{\de}=sk-z{\de}$ is the residue mod ${\de}$ of $sk$.
\end{lemma}

\begin{theorem}\label{T:powerassociative}
    Let $\alpha, b\in \ring$ and ${\de}\in \N$.  Suppose neither $\alpha$ nor $b$ are zero-divisors and $1\leq k <{\de}$.  Consider the magma generated by the monomial
    $$
    \scz=\alpha t^k
    $$
    in $\mathbb{S}_{t^\de-b}$. 
    If $\scz$ is power associative then
\begin{equation}\label{E:associativityrelation}
    \alpha b = \sigma^{\de}(\alpha)\sigma^k(b).
    \end{equation}
    Conversely, if \eqref{E:associativityrelation} holds, then  for every $s\in \mathbb{N}$, the expression $(\alpha t^k)^s$ is well-defined in $\mathbb{S}_{t^\de-b}$ 
    (that is, $\scz$ is power associative) and we have
    \begin{equation}\label{G(t)^s}
    \scz^s = N_{s}^{\sigma^k}(\alpha)\prod_{i=1}^{\lfloor \frac{sk}{{\de}}\rfloor} \sigma^{sk-i{\de}}(b)t^{[sk]_{\de}}
    \end{equation}
    where $\lfloor \frac{sk}{\de}\rfloor=j$ when $j{\de}\leq sk<(j+1){\de}$ and $[sk]_{\de} = sk-\lfloor \frac{sk}{\de}\rfloor \de$ denotes the residue of $sk$ mod $\de$ in the interval $[0,\de-1]$.
\end{theorem}

\begin{proof}
We first prove that if \eqref{E:associativityrelation} holds, then $\scz$ generates a monoid (that is, all powers are well-defined).
    Write $L(\scz,s)$ for the left-nested product
    $$
    \scz(\scz(\scz( \cdots (\scz))))
    $$
    of $\scz$ with itself $s$ times.
    Our first step is to show that this is equal to \eqref{G(t)^s} by induction on $s$.  The case $s=1$ is given.  Assume for some $s\geq 1$ that $L(\scz,s)$ is given by \eqref{G(t)^s}.  We wish to show that
    \begin{equation}\label{E:Lzs+1}
    L(\scz,s+1)= N_{s+1}^{\sigma^k}(\alpha)\prod_{i=1}^{\lfloor \frac{(s+1)k}{\de}\rfloor} \sigma^{(s+1)k-i\de}(b)t^{[(s+1)k]_{\de}}.
    \end{equation}
    We compute
    \begin{align*}
    L(\scz,s+1) &= \scz \cdot L(\scz,s)\\
    &=\alpha t^k \cdot N_{s}^{\sigma^k}(\alpha)\prod_{i=1}^{\lfloor \frac{sk}{\de}\rfloor} \sigma^{sk-i\de}(b)t^{[sk]_{\de}} \\
    &= \alpha \cdot \sigma^k\left( N_{s}^{\sigma^k}(\alpha)\right) \sigma^k\left(\prod_{i=1}^{\lfloor \frac{sk}{\de}\rfloor} \sigma^{sk-i\de}(b)\right)\cdot t^kt^{[sk]_{\de}}\\
    &= N_{s+1}^{\sigma^k}(\alpha) \prod_{i=1}^{\lfloor \frac{sk}{\de}\rfloor} \sigma^{(s+1)k-i\de}(b)\cdot t^kt^{[sk]_{\de}}.
    \end{align*}
    Note that since both $k$ and $[sk]_{\de}$ are less than $\de$, we simply have
    $$
    [(s+1)k]_{\de} = \begin{cases}
        k+[sk]_{\de} & \text{if $k+[sk]_{\de}<\de$, that is, if $\lfloor \frac{(s+1)k}{\de}\rfloor=\lfloor \frac{sk}{\de}\rfloor$};\\
        k+[sk]_{\de}-\de & \text{if $k+[sk]_{\de} > \de$, that is, if $\lfloor \frac{(s+1)k}{\de}\rfloor=\lfloor \frac{sk}{\de}\rfloor+1$.}
    \end{cases}
    $$
    In the former case we are done.  In the latter case, we apply the relation
    $$
    t^{k+[sk]_{\de}} = t^{k+[sk]_{\de}-\de}t^{\de}= t^{[(s+1)k]_{\de}}b=\sigma^{[(s+1)k]_{\de}}(b)t^{[(s+1)k]_{\de}},
    $$
    which gives the desired expression for $L(\scz,s+1),$  as required.

    We now assume that \eqref{E:associativityrelation} holds, and show that this implies power-associativity.  By induction it suffices to prove that
    for all $\ell,s\in \mathbb{N}$ we have
    $$
    L(\scz,s)L(\scz,\ell) = L(\scz,s+\ell).
    $$
This holds for $s=1$ and every $\ell\in \mathbb{N}$.  We compute
\begin{align*}
L&(\scz,s)\; L(\scz,\ell)=\\
&=N_{s}^{\sigma^k}(\alpha)\prod_{i=1}^{\lfloor \frac{s k}{\de}\rfloor} \sigma^{s k-i\de}(b)t^{[s k]_{\de}}\cdot N_{\ell}^{\sigma^k}(\alpha)\prod_{i=1}^{\lfloor \frac{\ell k}{\de}\rfloor} \sigma^{\ell k-i\de}(b)t^{[\ell k]_{\de}}\\
&= N_{s}^{\sigma^k}(\alpha)\prod_{i=1}^{\lfloor \frac{s k}{\de}\rfloor} \sigma^{s k-i\de}(b)\cdot
\sigma^{[s k]_{\de}}\left(N_{\ell}^{\sigma^k}(\alpha)\prod_{i=1}^{\lfloor \frac{\ell k}{\de}\rfloor} \sigma^{\ell k-i\de}(b)\right) \cdot t^{[s k]_{\de} + [\ell k]_{\de}}
\end{align*}
By Lemma~\ref{L:keystep}, our hypothesis \eqref{E:associativityrelation} implies
$$
\sigma^{[sk]_\de}(N_{\ell}^{\sigma^k}(\alpha))\prod_{j=1}^{\lfloor \frac{s k}{\de}\rfloor}\sigma^{sk-j\de}(b) = \sigma^{sk}(N_{\ell}^{\sigma^k}(\alpha))\prod_{j=1}^{\lfloor \frac{s k}{\de}\rfloor}\sigma^{sk-j\de+\ell k}(b).
$$
Thus
\begin{align*}
  L&(\scz,s)\; L(\scz,\ell)=\\
  &= N_{s}^{\sigma^k}(\alpha) \left(\prod_{i=1}^{\lfloor \frac{s k}{\de}\rfloor} \sigma^{s k-i\de}(b)\cdot
\sigma^{[s k]_{\de}}(N_{\ell}^{\sigma^k}(\alpha)) \right) \sigma^{[sk]_{\de}}\left(\prod_{i=1}^{\lfloor \frac{\ell k}{\de}\rfloor} \sigma^{\ell k-i\de}(b)\right) \cdot t^{[s k]_{\de} + [\ell k]_{\de}}\\
&= N_{s}^{\sigma^k}(\alpha)
\left( \sigma^{sk}(N_{\ell}^{\sigma^k}(\alpha))\prod_{j=1}^{\lfloor \frac{s k}{\de}\rfloor}\sigma^{sk-j\de+\ell k}(b)
\right)
\prod_{i=1}^{\lfloor \frac{\ell k}{\de}\rfloor} \sigma^{[sk]_\de+\ell k-i\de}(b) \cdot t^{[s k]_{\de} + [\ell k]_{\de}}.
\end{align*}
Since $[sk]_{\de}=sk-\lfloor \frac{sk}{\de}\rfloor \de$, we may reindex the last product to yield
\begin{align*}
&= N_{s+\ell}^{\sigma^k}(\alpha)
\prod_{j=1}^{\lfloor \frac{s k}{\de}\rfloor}\sigma^{sk-j\de +\ell k}(b)
\prod_{j=\lfloor \frac{s k}{\de}\rfloor+1}^{\lfloor \frac{s k}{\de}\rfloor+\lfloor \frac{\ell k}{\de}\rfloor}\sigma^{sk-j\de+\ell k}(b)  \cdot t^{[s k]_{\de} + [\ell k]_{\de}}\\
&= N_{s+\ell}^{\sigma^k}(\alpha)\prod_{j=1}^{\lfloor \frac{s k}{\de}\rfloor+\lfloor \frac{\ell k}{\de}\rfloor} \sigma^{(s+\ell)k-j\de }(b)  \cdot t^{[s k]_{\de} + [\ell k]_{\de}}.
\end{align*}
If $[sk]_{\de}+[\ell k]_{\de} = [(s+\ell)k]_{\de}$ then this is simply
$$
N_{s+\ell}^{\sigma^k}(\alpha)\prod_{j=1}^{\lfloor \frac{(s+\ell)k}{\de}\rfloor } \sigma^{(s+\ell)k-j\de}(b)  \cdot t^{[(s+\ell) k]_{\de}} = L(\scz,s+\ell).
$$
Otherwise, we must have
$$
\lfloor \frac{(s+\ell)k}{\de}\rfloor  = \lfloor \frac{s k}{\de}\rfloor+\lfloor \frac{\ell k}{\de}\rfloor+1,
$$
whence $
[sk]_{\de}+[\ell k]_{\de} = [(s+\ell)k]_{\de} + \de.
$
Since $[(s+\ell) k]_{\de}=(s+\ell)k-\lfloor \frac{(s+\ell)k}{\de}\rfloor \de$ we have
$$
t^{[s k]_{\de} + [\ell k]_{\de}}=t^{[(s+\ell) k]_{\de}+\de}= \sigma^{(s+\ell)k-\lfloor \frac{(s+\ell)k}{\de}\rfloor \de}(b)t^{[(s+\ell) k]_{\de}},
$$
which again yields $L(\scz,s)L(\scz,\ell)=L(\scz,s+\ell).$  Thus the desired equality holds for all $s,\ell \in \mathbb{N}$, and $\scz$ is power associative.

We now prove the condition is necessary.  Suppose $\scz$ is power-associative.  Then  $\scz\scz^r=\scz^r\scz$, for $r>0$ the least integer such that $rk\geq \de$.  We show that then  \eqref{E:associativityrelation} holds.  We begin by noting that $\scz^r$ is well-defined, since for all $s,\ell\geq 1$ such that $s+\ell=r$,
$$
\scz^s=N_s^{\sigma^k}(\alpha)t^{sk}, \quad \text{and} \quad \scz^\ell = N_\ell^{\sigma^k}(\alpha)t^{\ell k},
$$
by the associativity of powers less than $\de$ in $\Sone{b}$.  Thus we readily compute as above that
$$
\scz^r=\scz^s\scz^\ell= N_s^{\sigma^k}(\alpha)\sigma^{sk}(N_\ell^{\sigma^k}(\alpha))t^{(s+\ell)k}=N_r^{\sigma^k}(\alpha)\sigma^{rk-\de}(b)t^{rk-\de},
$$
independently of the choice of association, and this is equal to $L(\scz,r)$.

Now by our choice of $r$ (and $k$), we have $\de<(r+1)k < 2\de$ so that $\lfloor (r+1)/\de \rfloor=1$ and $[(r+1)k]_{\de}=(r+1)k-\de$.  Thus \eqref{E:Lzs+1} simplifies to
$$
\scz \scz^r=L(\scz,r+1)=N_{r+1}^{\sigma^k}(\alpha)\sigma^{(r+1)k-\de}(b)t^{(r+1)k-\de}
$$
whereas
$$  \scz^r \scz=N_r^{\sigma^k}(\alpha) \sigma^{rk-\de}(b)t^{rk-\de} \alpha t^k
= N_r^{\sigma^k}(\alpha) \sigma^{rk-\de}(\alpha)\sigma^{rk-\de}(b)t^{(r+1)k-\de}.
$$
If these are equal then we may equate these coefficients of $t^{(r+1)k-\de}$. If we divide both expressions by $N_r^{\sigma^k}(\alpha)$ and simplify, we obtain the equality
$$
\sigma^{rk}(\alpha)\sigma^{rk-\de}(\sigma^k(b))=\sigma^{rk-\de}(\alpha)\sigma^{rk-\de}(b),
$$
which, upon applying the automorphism $\sigma^{\de-rk}$ to both sides, is simply \eqref{E:associativityrelation}.
\end{proof}

\end{document}
